% Dirichlet Round 06 German TeX
% Work: G. Lejeune Dirichlet, Werke II, pp. 69--88.
\documentclass[11pt]{article}
\usepackage[a4paper,margin=1in]{geometry}
\usepackage[T1]{fontenc}
\usepackage[utf8]{inputenc}
\usepackage{lmodern}
\usepackage{amsmath,amssymb}
\usepackage{microtype}
\usepackage{hyperref}
\hypersetup{colorlinks=true,linkcolor=black,urlcolor=black,citecolor=black}
\setlength{\parindent}{1.2em}
\setlength{\parskip}{0.25em}
\newcommand{\dd}{\,d}
\newcommand{\half}{\frac12}
\begin{document}

\begin{center}
{\Large G. Lejeune Dirichlet's Werke. II.}\par\vspace{1.2em}
{\Large VII.}\par\vspace{0.5em}
{\LARGE \textbf{Über einen neuen Ausdruck}}\par
{\LARGE \textbf{zur Bestimmung der Dichtigkeit}}\par
{\LARGE \textbf{einer unendlich dünnen Kugelschale,}}\par
{\LARGE \textbf{wenn der Werth des Potentials derselben}}\par
{\LARGE \textbf{in jedem Punkte ihrer Oberfläche gegeben ist.}}\par\vspace{1em}
{\large Von G. Lejeune Dirichlet.}\par\vspace{0.5em}
{\small Abhandlungen der Königlich Preussischen Akademie der Wissenschaften von 1850, S. 99--116.}
\end{center}

\begin{center}
{\large \textbf{Über einen neuen Ausdruck zur Bestimmung}}\par
{\large \textbf{der Dichtigkeit einer unendlich dünnen Kugelschale,}}\par
{\large \textbf{wenn der Werth des Potentials derselben}}\par
{\large \textbf{in jedem Punkte ihrer Oberfläche gegeben ist.}}\par
{\small [Gelesen in der Akademie der Wissenschaften am 28. November 1850.\footnotemark]}
\end{center}
\footnotetext{Im Akademie-Bericht von 1850 wird auf S. 464 darüber mitgetheilt: „Hr. Dirichlet hielt einen Vortrag über einen neuen Ausdruck der Massen-Vertheilung auf einer Kugelfläche, wenn das Potential in jedem Punkte der Fläche einen beliebig gegebenen Werth erhalten soll.“ K.}

Nach einem schönen, von Gauss zuerst aufgestellten und bewiesenen Satze\footnote{Allgemeine Lehrsätze in Beziehung auf die im verkehrten Verhältnisse des Quadrats der Entfernung wirkenden Anziehungs- und Abstossungs-Kräfte von C. F. Gauss.} lässt sich jede Fläche mit einer unendlich dünnen Massenschicht belegen, deren Potential in jedem Punkte der Fläche einen beliebig gegebenen Werth erhält, vorausgesetzt dass dieser Werth im ganzen Umfange der Fläche sich nach der Stetigkeit ändere. Die wirkliche Ausmittelung aber der dies leistenden Massenvertheilung ist im gegenwärtigen Zustande der Wissenschaft nur für einige besondere Flächen ausführbar, zu welchen, wie Gauss schon bemerkt hat, die ganze Kugelfläche gehört.

Bezeichnet $R$ den Kugelradius, $r$ die Entfernung jedes Punktes im Raume vom Mittelpunkte, $\theta$ den Winkel zwischen $r$ und einer festen Geraden und $\varphi$ den Flächenwinkel zwischen der durch diese und $r$ gelegten Ebene und einer festen Ebene, so kann man den auf der ganzen Fläche gegebenen durch $\theta$ und $\varphi$ ausgedrückten Potentialwerth $V$ nach den bekannten Kugelfunctionen entwickeln. Es sei:
\[
  V=\sum X_n
\]
diese Entwickelung, wo sich das Summenzeichen von $n=0$ bis $n=\infty$ erstreckt, und:
\[
  \rho=\sum Y_n
\]
die ähnliche Entwickelung der zu bestimmenden Dichtigkeit $\rho$. Mit Hülfe dieser letzteren lässt sich das Potential $v$ für jeden Punkt im Raume darstellen.\footnote{Mécanique céleste, Livre III, No. 13.} Von den beiden Ausdrücken desselben, von welchen der eine im inneren, der andere im äusseren Raume gilt und jeder zu unserem Zwecke genügt, ist der erstere:
\[
 v=4\pi R\sum \frac{1}{2n+1}\left(\frac rR\right)^n Y_n.
\]
Da dieser Ausdruck bis $r=R$ gültig bleibt, für welchen Fall $v$ in $V$ übergeht, und dieselbe Function nur einer Entwickelung nach Kugelfunctionen fähig ist, so ergiebt die Vergleichung mit der ersten Reihe
\[
  Y_n=\frac{2n+1}{4\pi R}X_n,
\]
und daher:
\[
  \rho=\frac{1}{4\pi R}\sum(2n+1)X_n.
\]

In einer früheren Abhandlung\footnote{Sur les séries dont le terme général dépend de deux angles, et qui servent à exprimer des fonctions arbitraires entre des limites données. Crelle's Journal Band XVII. Bd. I, S. 283 dieser Ausgabe von G. Lejeune Dirichlet's Werken. K.} ist gezeigt worden, dass jede für die ganze Kugelfläche, d. h. von $\theta=0$, $\varphi=0$ bis $\theta=\pi$, $\varphi=2\pi$, beliebig gegebene Function, wenn sie nur nirgends unendlich wird, convergirend nach Kugelfunctionen entwickelt werden kann. Die Convergenz der Reihe für $V$, welche den Ausgangspunkt der eben entwickelten Lösung bildet, ist daher unzweifelhaft. Anders verhält es sich aber mit der für die Dichtigkeit $\rho$ angenommenen und dann durch Vergleichung mit jener bestimmten Reihe $\sum Y_n$. Da es mit der Existenz einer völlig bestimmten Massenvertheilung, welcher der gegebene Potentialwerth $V$ entspricht, verträglich ist, dass die Dichtigkeit in einzelnen Punkten oder Linien unendlich werde, so bleibt es für alle solche Fälle völlig ungewiss, ob die Reihe an allen Stellen, wo die Dichtigkeit endlich bleibt, ihre Convergenz behält und die Dichtigkeit wirklich darstellt. Es schien mir nicht uninteressant, diese Frage einer näheren Erörterung zu unterwerfen, wozu meine frühere Abhandlung die nöthigen Hilfsmittel darbot, um so mehr als sich von dieser Untersuchung eine Vereinfachung des eben für $\rho$ gefundenen Ausdrucks erwarten liess. Dieser Ausdruck involvirt eine vierfache unendliche Operation, eine doppelte Summation und eine ebenfalls doppelte Integration. Dass die letztere nicht aus dem Ausdrucke entfernt werden kann, leuchtet ein, da die Dichtigkeit für jeden Punkt von sämmtlichen im ganzen Umfange der Fläche bis auf die Stetigkeit als beliebig vorauszusetzenden Potentialwerthen abhängen muss. Dagegen hat sich ergeben, dass die Dichtigkeit immer ohne Reihen durch ein doppeltes Integral dargestellt werden kann, welches sogar in vielen Fällen auf ein einfaches zurückführbar ist.

\section*{1.}

Setzt man:
\[
 V=f(\theta,\varphi),
\]
so hat man bekanntlich für das allgemeine Glied $X_n$:
\[
 X_n=\frac{2n+1}{4\pi}\iint f(\theta',\varphi')P_n(\cos\omega)\sin\theta'\dd\theta'\dd\varphi',
\]
wo:
\[
 \cos\omega=\cos\theta\cos\theta' + \sin\theta\sin\theta'\cos(\varphi-\varphi')
\]
gesetzt ist, $P_n(\cos\omega)$ den Coefficienten von $\alpha^n$ in dem entwickelten Radical:
\[
 \frac{1}{\sqrt{1-2\alpha\cos\omega+\alpha^2}}
\]
bezeichnet und sich die doppelte Integration von
\[
 \theta'=0,\quad \varphi'=0 \qquad\text{bis}\qquad \theta'=\pi,\quad \varphi'=2\pi
\]
erstreckt. Lässt man den Divisor $R$ weg oder, was dasselbe ist, setzt den Kugelradius der Einheit gleich, so wird das allgemeine Glied der $\rho$ darstellenden Reihe:
\[
 \frac{(2n+1)^2}{(4\pi)^2}\iint f(\theta',\varphi')P_n(\cos\omega)\sin\theta'\dd\theta'\dd\varphi'.
\]
Es wird offenbar genügen, diese Reihe für den Fall zu untersuchen, wo $\theta=0$ gesetzt wird, d. h. für den Pol $p$ der sphärischen Polarcoordinaten $\theta$, $\varphi$, da sich, wie in der früheren Abhandlung, das für den Punkt $p$ gefundene Resultat unmittelbar auf jeden anderen Punkt $m$ der Fläche übertragen lässt. Man hat dann:
\[
 \cos\omega=\cos\theta',
\]
und $P_n(\cos\omega)$ wird von $\varphi'$ unabhängig. Setzt man:
\[
 \frac{1}{2\pi}\int_0^{2\pi} f(\theta',\varphi')\dd\varphi'=F(\theta'),
\]
so dass also $F(\theta')$ den mittleren Werth des Potentials $V$ auf dem von $p$ als Mittelpunkt mit dem sphärischen Radius $\theta'$ beschriebenen Kreise bedeutet, und schreibt $\gamma$ statt $\theta'$ zur Uebereinstimmung mit der Bezeichnung in der früheren Abhandlung, auf die wir uns häufig zu beziehen haben, so wird das allgemeine Glied:
\[
 \frac{(2n+1)^2}{8\pi}\int_0^\pi F(\gamma)P_n(\cos\gamma)\sin\gamma\dd\gamma.
\]

Zerlegt man $(2n+1)^2$ in seine Bestandtheile $4n^2$, $4n$, $1$, so zerfällt das allgemeine Glied in drei andere und folglich die Reihe selbst in drei Partialreihen. Da die Convergenz der zweiten und dritten dieser Reihen schon in der früheren Abhandlung gezeigt worden ist, so haben wir es nur mit der ersten zu thun, deren allgemeines Glied:
\[
 \frac{1}{2\pi}n^2\int_0^\pi F(\gamma)P_n(\cos\gamma)\sin\gamma\dd\gamma.
\]
Nimmt man die Summe der $n+1$ ersten Glieder, drückt $P_n(\cos\gamma)$ nach Gleichung (3) der früheren Abhandlung\footnote{Bd. I, S. 291 dieser Ausgabe von G. Lejeune Dirichlet's Werken. K.} durch ein bestimmtes Integral aus und kehrt die Ordnung der beiden Integrationen um, so erhält man:
\[
 \frac{1}{\pi^2}\int_0^\pi (\cos\psi+4\cos2\psi+\cdots+n^2\cos n\psi)\Pi(\psi)\dd\psi,
\]
wo wie früher:
\[
 \Pi(\psi)=\sin\frac{\psi}{2}\int_0^\psi F(\gamma)\sin\gamma\frac{\dd\gamma}{\Delta}
 +\cos\frac{\psi}{2}\int_\psi^\pi F(\gamma)\sin\gamma\frac{\dd\gamma}{E},
\]
\[
 \Delta=\sqrt{2(\cos\gamma-\cos\psi)},\qquad E=\sqrt{2(\cos\psi-\cos\gamma)}
\]
gesetzt ist. Da die in $\Pi(\psi)$ enthaltenen Integrale (nach §. 3 der früheren Abhandlung) Functionen von $\psi$ sind, welche von $\psi=0$ bis $\psi=\pi$ stetig bleiben, so kommt dieselbe Eigenschaft auch der Function $\Pi(\psi)$ zu.

Unsere gegenwärtige Untersuchung erfordert überdies die Discussion des Differentialquotienten $\Pi'(\psi)$, zu dessen Bildung, wegen der in den Integralen enthaltenen Nenner $\Delta$ und $E$, eine vorgängige Umformung durch theilweise Integration erforderlich ist. Berücksichtigt man, dass $F(\gamma)$, nach seinem Ursprunge aus dem stetigen Potentialwerthe $f(\theta,\varphi)$, selbst stetig ist, so ergiebt diese doppelte Operation:
\begin{align*}
\Pi'(\psi)=&\,(F(0)-F(\pi))\sin\psi
+\frac12\cos\frac{\psi}{2}\int_0^\psi F'(\gamma)\Delta\dd\gamma
+\frac12\sin\frac{\psi}{2}\int_\psi^\pi F'(\gamma)E\dd\gamma \\
&+\sin\frac{\psi}{2}\sin\psi\int_0^\psi F'(\gamma)\frac{\dd\gamma}{\Delta}
+\cos\frac{\psi}{2}\sin\psi\int_\psi^\pi F'(\gamma)\frac{\dd\gamma}{E}.
\end{align*}
Wir machen nun die Annahme, dass $F'(\gamma)$ von $\gamma=0$ bis $\gamma=\pi$ überall endlich bleibt. Alsdann werden sämmtliche Bestandtheile von $\Pi'(\psi)$, und also auch $\Pi'(\psi)$ selbst, von $\psi=0$ bis $\psi=\pi$ endlich und stetig sein. Für die drei ersten Glieder ist dies einleuchtend, und hinsichtlich der beiden letzten überzeugt man sich durch Betrachtungen, welche den oben citirten ganz ähnlich sind, dass die in ihnen enthaltenen Integrale diese doppelte Eigenschaft so lange behalten, als im ersten $\psi<\pi$ und im zweiten $\psi>0$ vorausgesetzt wird. Für $\psi=\pi$ und $\psi=0$ können zwar diese Integrale beziehungsweise unendlich grosse Werthe erhalten, die aber, wie leicht zu sehen, resp. die Ordnung
\[
 \log\left(\frac1{\cos\frac{\psi}{2}}\right)\qquad\text{und}\qquad
 \log\left(\frac1{\sin\frac{\psi}{2}}\right)
\]
nicht überschreiten können, so dass die Glieder selbst Null werden.

Für das Folgende ist noch die Kenntniss der zweiten Derivirten $\Pi''(\psi)$ für den besonderen Werth $\psi=0$ erforderlich. Bemerkt man, dass $\Pi'(0)=0$ ist, so lässt sich der gesuchte Werth am leichtesten dadurch finden, dass man $\psi$ in dem Quotienten $\frac1\psi\Pi'(\psi)$ unendlich klein werden lässt. Man erhält so:
\[
 \Pi''(0)=F(0)-F(\pi)+\frac12\int_0^\pi F'(\gamma)\sin\frac{\gamma}{2}\dd\gamma
 +\frac12\int_0^\pi \frac{F'(\gamma)}{\sin\frac{\gamma}{2}}\dd\gamma,
\]
oder, wenn man das dritte Glied theilweise integrirt:
\[
 \Pi''(0)=F(0)-\frac12 F(\pi)-\frac14\int_0^\pi F(\gamma)\cos\frac{\gamma}{2}\dd\gamma
 +\frac12\int_0^\pi \frac{F'(\gamma)}{\sin\frac{\gamma}{2}}\dd\gamma.
\]
Findet nun ausser der vorhin angenommenen Endlichkeit von $F'(\gamma)$ und der daraus folgenden Stetigkeit von $\Pi'(\psi)$ noch ein bestimmter endlicher Werth für $\Pi''(0)$ statt, oder, was dasselbe ist, hat das zweite Integral einen solchen Werth, so wird unsere Reihe immer convergiren und die Summe derselben leicht anzugeben sein. Setzt man nämlich zur Abkürzung:
\[
 X(\psi)=\frac{\sin(2n+1)\frac{\psi}{2}}{2\sin\frac{\psi}{2}}=\frac12+\cos\psi+\cos2\psi+\cdots+\cos n\psi,
\]
so wird obiger Ausdruck für die Summe der $n+1$ ersten Glieder:
\[
 -\frac1{\pi^2}\int_0^\pi \Pi(\psi)X''(\psi)\dd\psi.
\]
Da das bei zweimaliger theilweiser Integration heraustretende Glied:
\[
 \Pi(\psi)X'(\psi)-\Pi'(\psi)X(\psi)
\]
stetig ist und an beiden Grenzen verschwindet, so kommt durch diese Operation:
\[
 -\frac1{\pi^2}\int_0^\pi \Pi''(\psi)\frac{\sin(2n+1)\frac{\psi}{2}}{2\sin\frac{\psi}{2}}\dd\psi,
\]
und dieser Ausdruck nähert sich nach einem bekannten Satze, selbst dann, wenn $\Pi''(\psi)$ für einen oder mehrere von Null verschiedene Werthe von $\psi$ unendlich wird, bei wachsendem $n$ der Grenze:
\begin{align*}
 -\frac1{2\pi}\Pi''(0)=&-\frac1{2\pi}F(0)+\frac1{4\pi}F(\pi)
 +\frac1{8\pi}\int_0^\pi F(\gamma)\cos\frac{\gamma}{2}\dd\gamma \\
&-\frac1{4\pi}\int_0^\pi \frac{F'(\gamma)}{\sin\frac{\gamma}{2}}\dd\gamma.
\end{align*}
Addirt man zu diesem Werthe die der beiden anderen Reihen, wie sie in der früheren Abhandlung gefunden worden, so erhält man für die Dichtigkeit in $p$:
\[
 \rho=\frac1{4\pi}\left[F(\pi)-\int_0^\pi \frac{F'(\gamma)}{\sin\frac{\gamma}{2}}\dd\gamma\right].
\]

Es bliebe nun noch zu untersuchen, wie es sich mit der die Dichtigkeit ausdrückenden Reihe rücksichtlich ihrer Convergenz in dem Falle verhält, wo zwar das Integral
\[
 \int_0^\pi \frac{F'(\gamma)}{\sin\frac{\gamma}{2}}\dd\gamma
\]
und also auch der eben gefundene Ausdruck für $\rho$ einen bestimmten endlichen Werth behält, die Function $F'(\gamma)$ aber für einen oder mehrere von Null verschiedene Werthe von $\gamma$ unendlich wird. Die Bedingung für die Convergenz besteht alsdann nach Obigem lediglich darin, dass die aus $F'(\gamma)$ abgeleitete Function $\Pi''(\psi)$ von $\psi=0$ bis $\psi=\pi$ endlich und stetig bleibe. Eine genauere Betrachtung der Bildungsweise von $\Pi''(\psi)$ ergiebt nun, dass, wenn das Unendlichwerden von $F'(\gamma)$ so erfolgt, dass für jeden Werth $c$, für den $F'(\gamma)$ einen unendlich grossen Werth erhält, $\sqrt{\varepsilon}F'(c\pm\varepsilon)$ für ein unendlich kleines $\varepsilon$ selbst unendlich klein wird, die Continuität von $\Pi''(\psi)$ ebenso stattfindet wie in dem vorhin untersuchten Falle einer überall endlichen Function $F'(\gamma)$ und mit ihr die Convergenz der Reihe, welche die Dichtigkeit darstellt, dass hingegen im Allgemeinen Divergenz eintritt, wenn die eben ausgesprochene Bedingung nicht mehr erfüllt ist. Obgleich die Begründung des erwähnten Resultates keine wesentlichen Schwierigkeiten darbietet, so erfordert sie doch andererseits zu viel Raum und gewährt zu wenig Interesse, um dieselbe hier durchzuführen. Es genügt für den sicheren Gebrauch der Reihe, durch das Obige die Fälle zu kennen, in denen allein die Convergenz der Reihe aufhören kann. Uebrigens wird sich weiter unten Gelegenheit finden, die in einem solchen Falle wirklich eintretende Divergenz an einem höchst einfachen Beispiele nachzuweisen.

\section*{2.}

Da die vorige Ableitung des für $\rho$ gefundenen endlichen Ausdrucks ihre Gültigkeit verliert, wenn die Reihe zu convergiren aufhört, so bedarf es noch eines auf alle Fälle anwendbaren Beweises dieses Ausdrucks.

Nach einem bekannten Satze nähert sich der Differentialquotient $\frac{\partial v}{\partial r}$, wenn darin $\theta$ und $\varphi$ als constant, $r$ aber der Einheit sich nähernd gedacht werden, zwei verschiedenen Grenzen, je nachdem dabei immerfort $r<1$ oder $r>1$ vorausgesetzt wird. Nennt man diese Grenzwerthe beziehungsweise $K$ und $L$, so wird die Dichtigkeit im entsprechenden Punkte der Fläche durch die Gleichung:
\[
 \rho=\frac1{4\pi}(K-L)
\]
bestimmt. Zwischen den Grenzen $K$, $L$ und dem Potentialwerth $V$ auf der Fläche findet ein einfacher Zusammenhang statt, so dass man es nur mit der Ausmittelung eines der Grenzwerthe zu thun hat. Sind nämlich $v$ und $v_1$ die Potentialwerthe für zwei Punkte auf demselben Radiusvector, deren Entfernungen vom Mittelpunkt, $r$ und $r_1$, der Gleichung
\[
 rr_1=1
\]
genügen, so hat man, wie leicht ersichtlich:
\[
 v_1=\frac1{r_1}v
\]
und folglich:
\[
 \frac{\partial v_1}{\partial r_1}=-\frac1{r_1^2}v+\frac1{r_1}\frac{\partial v}{\partial r}\frac{\partial r}{\partial r_1}
 =-\frac1{r_1^2}v-\frac1{r_1^3}\frac{\partial v}{\partial r}.
\]
Man sieht also, dass:
\[
 L=-V-K
\]
und:
\[
 \rho=\frac1{4\pi}(V+2K)
\]
ist. Zur Bestimmung von $K$ bedarf es nur des für den inneren Raum geltenden Ausdrucks von $v$:
\[
 v=\frac{1-r^2}{4\pi}\iint\frac{f(\theta',\varphi')\sin\theta'\dd\theta'\dd\varphi'}{(1-2r\cos\omega+r^2)^{3/2}},
\]
welcher leicht ohne Reihenentwickelung bewiesen wird. Es genügt dazu die Bemerkung, dass der Ausdruck, wenn in demselben $r$ sich seiner oberen Grenze $1$ nähert, nach einem bekannten vielfach behandelten Satze in $f(\theta,\varphi)=V$ übergeht, und dass derselbe andererseits der zuerst von Laplace für $v$ aufgestellten partiellen Differentialgleichung genügt, da:
\[
 \frac{1-r^2}{(1-2r\cos\omega+r^2)^{3/2}},
\]
als Function der Polarcoordinaten $r$, $\theta$, $\varphi$ betrachtet, ein particuliäres Integral dieser Gleichung darstellt. Für den Fall, wo $\theta=0$ ist, nimmt der Ausdruck, wenn wieder $\gamma$ statt $\theta'$ geschrieben und die frühere Bezeichnung beibehalten wird, die einfache Form an:
\[
 v=\frac12(1-r^2)\int_0^\pi \frac{F(\gamma)\sin\gamma\dd\gamma}{(1-2r\cos\gamma+r^2)^{3/2}}.
\]
Differentiirte man den Ausdruck in dieser Gestalt, so würde sich der Grenzwerth von $\frac{\partial v}{\partial r}$ schwer bestimmen lassen, und es ist zweckmässig, vorher eine theilweise Integration mit demselben vorzunehmen. Man erhält so:
\[
 2v=\left(\frac1r+1\right)F(0)-\left(\frac1r-1\right)F(\pi)
 +\left(\frac1r-r\right)\int_0^\pi\frac{F'(\gamma)\dd\gamma}{(1-2r\cos\gamma+r^2)^{1/2}},
\]
wo jetzt der Uebergang zur Grenze nach vorher geschehener Differentiation keine Schwierigkeiten mehr darbietet und das Resultat ergiebt:
\[
 K=-\frac12F(0)+\frac12F(\pi)-\frac12\int_0^\pi\frac{F'(\gamma)}{\sin\frac{\gamma}{2}}\dd\gamma.
\]
Hieraus folgt mit Berücksichtigung, dass für den Punkt $p$:
\[
 V=F(0)
\]
ist, die Gleichung:
\[
 \rho=\frac1{4\pi}\left[F(\pi)-\int_0^\pi\frac{F'(\gamma)}{\sin\frac{\gamma}{2}}\dd\gamma\right],
\]
welche mit der aus der Reihe abgeleiteten übereinstimmt.

\section*{3.}

Wird das Resultat, welches für den $\theta=0$ entsprechenden Punkt gefunden worden ist, auf einen beliebigen Punkt $m$ übertragen, so ergiebt sich zur Bestimmung der Dichtigkeit die folgende allgemeine Regel.

Man suche durch eine erste Integration für den von $m$ als Mittelpunkt mit dem beliebigen sphärischen Radius $\lambda$ auf der Fläche beschriebenen Kreis den mittleren Werth des gegebenen Potentials $V$. Bezeichnet man diesen mit $\varphi(\lambda)$, so wird die gesuchte Dichtigkeit $\rho$ durch die Gleichung gegeben:
\[
 \rho=\frac1{4\pi}\left[\varphi(\pi)-\int_0^\pi\frac{\varphi'(\lambda)}{\sin\frac{\lambda}{2}}\dd\lambda\right].
\]

Man kann noch bemerken, dass, nach der Bedeutung der Function $\varphi(\lambda)$, der besondere Werth $\varphi(\pi)$ den Potentialwerth für den Gegenpunkt des Punktes $m$ bezeichnet, und es bedarf kaum der Erwähnung, dass die Function $\varphi(\lambda)$ für jeden Punkt $m$ eine andere sein, oder mit anderen Worten, dass $\varphi(\lambda)$ ausser $\lambda$ noch die beiden Grössen enthalten wird, durch welche die Lage von $m$ auf der Fläche bestimmt wird.

Dagegen ist es vielleicht nicht überflüssig, einem Irrthum vorzubeugen, welcher bei unaufmerksamer Betrachtung des eben ausgesprochenen Resultats leicht entstehen kann. Man ist auf den ersten Blick wegen des unter dem Integralzeichen vorkommenden, an der unteren Grenze verschwindenden Divisors versucht zu vermuthen, dass das Integral nur für besondere Lagen des Punktes $m$ endlich bleibt, im Allgemeinen aber unendlich wird, während gerade das umgekehrte Verhältniss stattfindet. Zunächst ist in Folge der Stetigkeit von $\varphi(\lambda)$ leicht einzusehen, dass der Theil des Integrals, welcher sich von einem noch so kleinen positiven Werthe $\delta$ bis zur oberen Grenze $\pi$ erstreckt, immer bestimmt und endlich bleibt, wie oft auch $\varphi'(\lambda)$ selbst in diesem Intervalle unendlich werde. Dass aber auch für den von $0$ bis $\delta$ sich erstreckenden Theil des Integrals, singuläre Fälle ausgenommen, dasselbe gilt, hat seinen Grund darin, dass $\varphi'(\lambda)$ für kleine Werthe von $\lambda$ im Allgemeinen wie der Nenner von der ersten Ordnung ist. Stellt man den Potentialwerth $V$ als Function $\chi(\lambda,\psi)$ der sphärischen Polarcoordinaten $\lambda$, $\psi$ dar, deren Pol der Punkt $m$ ist, wo dann:
\[
 \varphi(\lambda)=\frac1{2\pi}\int_0^{2\pi}\chi(\lambda,\psi)\dd\psi
\]
wird, so tritt die erwähnte Eigenschaft nur deshalb nicht hervor, weil die Function $\chi(\lambda,\psi)$ keine beliebige ist sondern die besonderen Bedingungen erfüllen muss, in Bezug auf $\psi$ periodisch zu sein und für $\lambda=0$ von $\psi$ unabhängig zu werden. Man überzeugt sich dagegen sogleich von dem vorhin Behaupteten, wenn man die Kugelfläche auf eine beliebige Ebene projicirt, deren Punkte auf zwei rechtwinklige Axen der $t$ und $u$ bezogen sind, und dann $V$ für jeden Punkt durch die Coordinaten $t$, $u$ seiner Projection ausdrückt. Es hat diese Darstellungsweise allerdings den Uebelstand, dass man, um die ganze Fläche zu umfassen, zwei Functionen von $t$ und $u$ zu betrachten hat; dieser Uebelstand fällt aber für unseren Zweck weg, der nur die Berücksichtigung eines beliebig kleinen, den Punkt $m$ einschliessenden Flächenstücks erfordert. Nimmt man die an $m$ gelegte Tangentialebene zur Projectionsebene, $m$ zum Anfangspunkt der Coordinaten und setzt $V=\chi(t,u)$, so hat diese neue Function keine andere Bedingung als die der Stetigkeit zu erfüllen. Da nun offenbar:
\[
 t=\sin\lambda\cos\psi,\qquad u=\sin\lambda\sin\psi
\]
angenommen werden kann, so erhält $\varphi(\lambda)$ die Form:
\[
 \varphi(\lambda)=\frac1{2\pi}\int_0^{2\pi}\chi(\sin\lambda\cos\psi,\sin\lambda\sin\psi)\dd\psi,
\]
in welcher das Stattfinden der oben erwähnten Eigenschaft einleuchtet, wenn anders die Differentialquotienten von $\chi(t,u)$ bis zur zweiten Ordnung incl. für den Fall, wo gleichzeitig $t=0$ und $u=0$ ist, bestimmte endliche Werthe behalten. Uebrigens ist es für die Endlichkeit von $\rho$ nicht einmal erforderlich, dass $\varphi'(\lambda)$ für kleine Werthe der Veränderlichen von der ersten Ordnung sei; es genügt, dass die Ordnung von $\varphi'(\lambda)$ mit der irgend einer positiven Potenz von $\lambda$ übereinstimme.

Aehnlich wie mit dem Unendlichwerden von $\rho$ und der daraus folgenden Divergenz der Reihe verhält es sich mit dem Falle, wo die Reihe bei endlich bleibender Dichtigkeit zu convergiren aufhört. Dieser Fall hat sogar noch weniger Umfang als der eben besprochene, wie man sich bei näherer Erwägung der Bedingungen, von denen er abhängt, sehr leicht überzeugen wird.

\section*{4.}

Um das wirkliche Stattfinden der Divergenz der Reihe für $\rho$ in den oben bezeichneten Ausnahmefällen durch ein einfaches Beispiel zu erläutern, machen wir die Voraussetzung, dass der Potentialwerth $V$ den Winkel $\varphi$ nicht enthält und durch $f(\theta)$ dargestellt wird. Es tritt dann bekanntlich eine einfachere Form für die Kugelfunctionreihen ein, und man hat:
\[
 V=\frac12\sum(2n+1)A_nP_n(\cos\theta),\qquad
 \rho=\frac1{8\pi}\sum(2n+1)^2A_nP_n(\cos\theta),
\]
wo der allgemeine Coefficient $A_n$ durch die Gleichung:
\[
 A_n=\int_0^\pi f(\theta)P_n(\cos\theta)\sin\theta\dd\theta
\]
gegeben wird. Setzt man:
\[
 f(\theta)=\sqrt{\cos\theta}\quad\text{so lange }0<\theta<\frac12\pi,
\qquad
 f(\theta)=0\quad\text{wenn }\theta\text{ zwischen }\frac12\pi\text{ und }\pi\text{ liegt},
\]
so wird der Forderung der Continuität genügt, und man hat das Integral:
\[
 A_n=\int_0^{\frac12\pi}P_n(\cos\theta)\sqrt{\cos\theta}\sin\theta\dd\theta
\]
zu bestimmen. Da die Ausmittelung desselben mit Hülfe der bekannten Ausdrücke für $P_n(\cos\theta)$ das Resultat nicht unmittelbar in der einfachsten Form ergiebt, so ist der folgende Weg vorzuziehen.

In Folge der Gleichung:
\[
 \sum P_n(\cos\theta)\alpha^n=\frac1{\sqrt{1-2\alpha\cos\theta+\alpha^2}}
\]
ist das gesuchte Integral der Coefficient von $\alpha^n$ in dem entwickelten Ausdrucke:
\[
 \int_0^{\frac12\pi}\frac{\sqrt{\cos\theta}\sin\theta\dd\theta}{\sqrt{1-2\alpha\cos\theta+\alpha^2}},
\]
in welchem der ächte Bruch $\alpha$ als positiv betrachtet werden kann. Durch die Substitution $t=\sqrt{\cos\theta}$ und Ausführung der Integration erhält man:
\[
 2\int_0^1\frac{t^2\dd t}{\sqrt{1-2\alpha t^2+\alpha^2}}
 =\sqrt{\frac1{8\alpha^3}}\left((1+\alpha^2)\arcsin\sqrt{\frac{2\alpha}{1+\alpha^2}}-(1-\alpha)\sqrt{2\alpha}\right).
\]
Setzt man zur Erleichterung der Entwickelung:
\[
 z=\arcsin\sqrt{\frac{2\alpha}{1+\alpha^2}}
\]
und differentiirt, so kommt:
\[
 \frac{dz}{d\alpha}=\frac{1+\alpha}{1+\alpha^2}\sqrt{\frac1{2\alpha}}
 =\frac1{\sqrt2}(\alpha^{-\frac12}+\alpha^{\frac12}-\alpha^{\frac32}-\alpha^{\frac52}+\cdots),
\]
und wenn man integrirt und bemerkt, dass $z$ mit $\alpha$ verschwindet:
\[
 \arcsin\sqrt{\frac{2\alpha}{1+\alpha^2}}
 =\sqrt2\left(\alpha^{\frac12}+\frac13\alpha^{\frac32}-\frac15\alpha^{\frac52}-\frac17\alpha^{\frac72}+\cdots\right).
\]
Setzt man ein, so erhält man für die gesuchte Entwickelung:
\[
 -\frac12\left(\left(\frac1{-1}-\frac13\right)-\left(\frac11-\frac15\right)\alpha-
 \left(\frac13-\frac17\right)\alpha^2+\left(\frac15-\frac19\right)\alpha^3+\cdots\right),
\]
wo die Vorzeichen innerhalb der Klammer die viergliedrige Periode
\[
 +\; -\; -\; +
\]
bilden, welche durch $(-1)^{\frac12 n(n+1)}$ dargestellt werden kann. Es ist also:
\[
 A_n=-(-1)^{\frac12 n(n+1)}\frac{2}{(2n-1)(2n+3)},
\]
und obige Reihen erhalten die Form:
\[
 V=-\sum (-1)^{\frac12 n(n+1)}\frac{2n+1}{(2n-1)(2n+3)}P_n(\cos\theta),
\]
\[
 \rho=-\frac1{4\pi}\sum (-1)^{\frac12 n(n+1)}\frac{(2n+1)^2}{(2n-1)(2n+3)}P_n(\cos\theta).
\]
Nach dem im vorigen Artikel Bemerkten sieht man in unserem Falle ohne Schwierigkeit, dass der Differentialquotient $\varphi'(\lambda)$ nur dann, und zwar für $\lambda=\frac12\pi$, so unendlich wird, dass die am Ende von Art. 1 ausgesprochene Bedingung nicht erfüllt ist, wenn der Punkt $m$ dem Werthe $\theta=0$ oder $\theta=\pi$ entspricht, in welchen Fällen jedoch $\varphi'(\lambda)$ für kleine Werthe von $\lambda$ die Ordnung $\lambda$ und also $\rho$ einen bestimmten endlichen Werth behält, so wie auch dass die Dichtigkeit nur am Aequator oder für $\theta=\frac12\pi$ unendlich wird, indem alsdann $\varphi'(\lambda)$ für ein kleines $\lambda$ von der Ordnung $\lambda^{-1}$ ist. In diesen drei Fällen wird also die Reihe zu convergiren aufhören, wovon man sich sogleich überzeugt, wenn man berücksichtigt, dass $P_n(\cos\theta)$ für $\theta=0$ und $\theta=\pi$ beziehungsweise die Werthe $1$ und $(-1)^n$ hat, für $\theta=\frac12\pi$ aber, wenn $n$ ungerade ist, verschwindet und, wenn $n$ gerade ist, dem Ausdrucke:
\[
 \frac{1\cdot3\cdot5\cdots(n-1)}{2\cdot4\cdot6\cdots n}(-1)^{\frac n2}
\]
gleich wird, welcher bekanntlich für grosse Werthe von $n$ von der Ordnung $n^{-\frac12}$ ist.

In den beiden ersten Fällen, wo $\rho$ einen bestimmten endlichen Werth behält, hat die Reihe den Charakter einer oscillirenden, indem ihre Glieder, von denen unter je vier auf einander folgenden zwei das positive und zwei das negative Zeichen haben, sich der Einheit als Grenze nähern, während im dritten Falle alle Glieder dasselbe Vorzeichen erhalten und eine unendlich grosse Summe ergeben.

\section*{5.}

Obgleich die Methode, durch welche wir vorhin den Coefficienten $A_n$ bestimmt haben, nicht mehr anwendbar ist, wenn man, unter Beibehaltung der Voraussetzung $f(\theta)=0$ für den zweiten Theil des Intervalls, im ersten $f(\theta)=\cos^k\theta$ annimmt, wo $k$ eine beliebige positive Constante bezeichnet, so lässt die höchst einfache Form des für den besonderen Werth $k=\frac12$ gefundenen Resultats etwas Aehnliches für den allgemeineren Fall erwarten. Da der einfachste Ausdruck von $A_n$ jedoch etwas versteckt ist, so wollen wir uns einen Augenblick bei dessen Ausmittelung aufhalten. Poisson, welcher in einer seiner Abhandlungen\footnote{Connaissance des temps pour l'an 1829.} die eben definirte Function als Beispiel der Entwickelung nach Kugelfunctionen gewählt hat, lässt dem Coefficienten $A_n$ die Form:
\[
 A_n=\frac{1\cdot3\cdots(2n-1)}{1\cdot2\cdots n}\left(\frac1{k+n+1}-\frac{n(n-1)}{2(2n-1)}\frac1{k+n-1}
 +\frac{n(n-1)(n-2)(n-3)}{2\cdot4(2n-1)(2n-3)}\frac1{k+n-3}-\cdots\right),
\]
in welcher sich derselbe unmittelbar darstellt, wenn man bei der Entwickelung den gewöhnlichen Ausdruck von $P_n(\cos\theta)$ zu Grunde legt, ohne die grosse Vereinfachung zu bemerken, deren diese Form fähig ist, und welche man in der That nicht leicht vermuthet, wenn man nicht durch einen besonderen Fall, wie der obige, darauf aufmerksam gemacht wird. Zu dem einfachsten Ausdruck für $A_n$ führt der folgende sehr kurze, wenn auch nicht elementare Weg.

Da nach §. 1 der früheren Abhandlung\footnote{Bd. I, S. 292--294 dieser Ausgabe von G. Lejeune Dirichlet's Werken. K.} die beiden Integrale:
\[
 \frac{2}{\pi}\int_0^\gamma \frac{\cos n\psi\cos\frac{\psi}{2}}{\sqrt{2(\cos\psi-\cos\gamma)}}\dd\psi,
\qquad
 -\frac{2}{\pi}\int_0^\gamma \frac{\sin n\psi\sin\frac{\psi}{2}}{\sqrt{2(\cos\psi-\cos\gamma)}}\dd\psi
\]
resp. den Coefficienten von $\alpha^n$ in den Entwickelungen von:
\[
 \frac12+\frac12\frac{1+\alpha}{\sqrt{1-2\alpha\cos\gamma+\alpha^2}},
\qquad
 -\frac12+\frac12\frac{1-\alpha}{\sqrt{1-2\alpha\cos\gamma+\alpha^2}}
\]
gleich sind, so erhält man durch Addition:
\[
 P_n(\cos\gamma)=\frac2\pi\int_0^\gamma\frac{\cos(n+\frac12)\psi}{\sqrt{2(\cos\psi-\cos\gamma)}}\dd\psi.
\]
Wird diese Gleichung mit $\cos^k\gamma\sin\gamma\dd\gamma$ multiplicirt, alsdann von $\gamma=0$ bis $\gamma=\frac12\pi$ integrirt, und die Ordnung der Integrationen auf der zweiten Seite umgekehrt, so ergiebt sich:
\[
 A_n=\frac2\pi\int_0^{\frac12\pi}\dd\psi\,\cos(n+\tfrac12)\psi
 \int_\psi^{\frac12\pi}\frac{\cos^k\gamma\sin\gamma}{\sqrt{2(\cos\psi-\cos\gamma)}}\dd\gamma,
\]
oder, wenn man für $\gamma$ durch die Gleichung $\cos\gamma=t\cos\psi$ eine neue Veränderliche $t$ einführt, wodurch sich die beiden Integrationen von einander trennen:
\[
 A_n=\frac{\sqrt2}{\pi}\int_0^1\frac{t^k}{\sqrt{1-t}}\dd t\cdot
 \int_0^{\frac12\pi}\cos^{k+\frac12}\psi\cos(n+\tfrac12)\psi\dd\psi.
\]
Nun ist nach einer bekannten Formel:
\[
 \int_0^{\frac12\pi}\cos^{p-1}\psi\cos q\psi\dd\psi
 =\frac{\Gamma(p)}{\Gamma\left(\frac12(1+p+q)\right)\Gamma\left(\frac12(1+p-q)\right)}\cdot\frac{\pi}{2^p},
\]
in welcher die Constante $p$ positiv sein muss und die Transcendente $\Gamma(l)$, wenn das Argument $l$ negativ wird, nach der Relation $\Gamma(l+1)=l\Gamma(l)$, oder, was dasselbe ist, nach der von Gauss gegebenen Definition zu verstehen ist, welche positive und negative Argumente gleichmässig umfasst. Vermittelst dieser Formel und der bekannten Darstellung eines Eulerschen Integrals der ersten Gattung durch drei der zweiten, wird unsere Gleichung:
\[
 A_n=\frac{\Gamma(k+1)}{\Gamma\left(\frac12(k+n+3)\right)\Gamma\left(\frac12(k-n+2)\right)}\cdot\frac{\sqrt\pi}{2^{k+1}},
\]
oder:
\[
 A_n=\frac{k(k-1)(k-2)\cdots(k-n+1)}{\frac12(k+n+1)\left(\frac12(k+n+1)-1\right)\cdots\left(\frac12(k+n+1)-n\right)}
 \cdot\frac{\Gamma(k-n+1)}{\Gamma\left(\frac12(k-n+1)\right)\Gamma\left(\frac12(k-n+2)\right)}\cdot\frac{\sqrt\pi}{2^{k+1}},
\]
und wenn man für den zweiten Factor seinen bekannten Werth $2^{k-n}\pi^{-\frac12}$ setzt:
\[
 A_n=\frac{k(k-1)(k-2)\cdots(k-n+1)}{(k+n+1)(k+n-1)(k+n-3)\cdots(k-n+1)},
\]
wo die Factoren im Nenner um zwei Einheiten abnehmen.

Es bedarf kaum der Erwähnung, dass dieser einfache Ausdruck mit der grössten Leichtigkeit verificirt werden kann. Man darf denselben z. B. nur in Partialbrüche zerlegen, um die Uebereinstimmung mit dem oben erwähnten zu erkennen. Will man letzteren zu dieser Verification nicht benutzen, so kann man sich dazu der Gleichung:
\[
 (n+1)A_{n+1,k}=(2n+1)A_{n,k+1}-nA_{n-1,k}
\]
bedienen, die unmittelbar aus der Relation folgt, welche zwischen je drei aufeinander folgenden der Entwickelungscoefficienten $P_n(\cos\gamma)$ stattfindet.

Um den gefundenen Ausdruck noch weiter zu vereinfachen, hat man zu unterscheiden, ob $n$ gerade oder ungerade ist, und erhält beziehungsweise:
\[
 A_n=\frac{k(k-2)\cdots(k-n+2)}{(k+1)(k+3)\cdots(k+n+1)},
\qquad
 A_n=\frac{(k-1)(k-3)\cdots(k-n+2)}{(k+2)(k+4)\cdots(k+n+1)}.
\]
Mit Hülfe dieser Ausdrücke und der bekannten Formeln, welche die genäherten Werthe von Facultäten sehr grosser Factorenanzahl ergeben, kann man sich leicht überzeugen, dass der oben speciell untersuchte Fall $k=\frac12$ hinsichtlich der bei $\theta=0$ und $\theta=\pi$ stattfindenden Divergenz der Reihe gerade der Grenzfall ist, und dass die Divergenz aufhört, sobald $k>\frac12$ ist. Anders verhält es sich dagegen mit dem Werthe $\theta=\frac12\pi$, an dieser Stelle besteht die Divergenz und der unendlich grosse Werth der Dichtigkeit so lange fort, als nicht $k>1$ angenommen wird.

\section*{6.}

Wir wollen nun noch von dem für die Dichtigkeit gefundenen endlichen Ausdruck eine Anwendung auf einen speciellen Fall machen. Wird wieder $V$ als eine blosse Function von $\theta$ betrachtet, die wir in die Form $f(\cos\theta)$ bringen können, so wird auch $\rho$ nur von der sphärischen Entfernung $\theta$ des Punktes $m$ vom Pol $p$ abhängen. Auf dem von $m$ als Mittelpunkt mit dem sphärischen Radius $\lambda$ beschriebenen Kreise ist dann in Folge der Grundformel der Trigonometrie:
\[
 V=f(\cos\theta\cos\lambda+\sin\theta\sin\lambda\cos\psi),
\]
wo $\psi$ den Winkel zwischen einem beliebigen Radius $\lambda$ und dem von $m$ nach $p$ gerichteten bezeichnet. Da dieser Ausdruck für $\psi$ und $-\psi$ denselben Werth hat, so kann man das den mittleren Werth $\varphi(\lambda)$ ausdrückende Integral von $\psi=0$ bis $\psi=\pi$ nehmen und dann verdoppeln. Man erhält so:
\[
 \varphi(\lambda)=\frac1\pi\int_0^\pi f(\cos\theta\cos\lambda+\sin\theta\sin\lambda\cos\psi)\dd\psi.
\]
So oft sich $\varphi'(\lambda)$ ohne Integralzeichen darstellen lässt, was namentlich der Fall ist, wenn $f(z)$ oder noch allgemeiner $f'(z)$ eine rationale Function von $z$ ist, mag nun die Form dieser Function für das ganze Intervall zwischen $z=-1$ und $z=+1$ dieselbe bleiben oder sich stellenweise so ändern, dass die Stetigkeit nicht verletzt wird, wird $\rho$ auf ein einfaches Integral zurückgeführt oder selbst ohne Integralzeichen ausgedrückt werden können.\footnote{Im allgemeinen Falle, wo $V$ beide Winkel $\theta$, $\varphi$ enthält, findet die Zurückführung von $\rho$ auf eine einfache Quadratur immer statt, wenn man, $x=\cos\theta$, $y=\sin\theta\cos\varphi$, $z=\sin\theta\sin\varphi$ setzend, $V$ durch eine Function $f(x,y,z)$ der Grössen $x,y,z$ darstellen kann, deren drei Differentialquotienten erster Ordnung rationale, ganze oder gebrochene Functionen von $x,y,z$ sind, wobei wieder die Form von $f(x,y,z)$ in verschiedenen Theilen der Fläche eine andere sein kann, was ein durch seinen grossen Umfang bemerkenswerthes Resultat ist.}

Wir behandeln unter den hierher gehörigen Fällen nur den, wo:
\[
 f(\cos\theta)=\cos\theta,
\]
so lange $\theta<\frac12\pi$, und:
\[
 f(\cos\theta)=0,
\]
wenn $\theta>\frac12\pi$ angenommen wird. Das Integral wird sich dann nur über den Theil des Intervalls zwischen $\psi=0$ und $\psi=\pi$ erstrecken, innerhalb dessen:
\[
 \cos\theta\cos\lambda+\sin\theta\sin\lambda\cos\psi
\]
positive Werthe erhält. Setzt man $\theta<\frac12\pi$ voraus, auf welchen Fall der, wo $\theta>\frac12\pi$ ist, leicht zurückgeführt wird, so findet man durch eine einfache Discussion des vorigen Ausdrucks, dass, so lange $\lambda$ unter $\frac12\pi-\theta$ liegt, das Integral von $\psi=0$ bis $\psi=\pi$ zu nehmen ist, dass für die Werthe von $\lambda$, welche zwischen $\frac12\pi-\theta$ und $\frac12\pi+\theta$ liegen, die Grenzen des Integrals $0$ und der durch die Gleichung:
\[
 \cos\theta\cos\lambda+\sin\theta\sin\lambda\cos\psi_1=0
\]
bestimmte, zwischen $0$ und $\pi$ liegende Winkel $\psi_1$ sind, und dass endlich für $\lambda>\frac12\pi+\theta$ das Integral verschwindet, indem der obige Ausdruck immerfort negativ ist. Die Function $\varphi(\lambda)$ ist also nach den drei eben unterschiedenen Intervallen:
\[
 \cos\theta\cos\lambda,
\qquad
 \frac1\pi\int_0^{\psi_1}(\cos\theta\cos\lambda+\sin\theta\sin\lambda\cos\psi)\dd\psi,
\qquad
 0,
\]
wo im zweiten Falle die Ausführung der Integration zur Abkürzung der Rechnung bis nach der Differentiation aufgeschoben ist. Bei dieser verschwindet das von der Veränderlichkeit der oberen Grenze $\psi_1$ herrührende Glied in Folge der $\psi_1$ bestimmenden Gleichung, und man erhält für $\varphi'(\lambda)$ die drei Ausdrücke:
\[
 -\cos\theta\sin\lambda,
\qquad
 \frac1\pi(-\psi_1\cos\theta\sin\lambda+\sin\psi_1\sin\theta\cos\lambda),
\qquad
 0,
\]
an welchen man sich im Vorbeigehen überzeugen kann, dass, wenigstens so lange $m$ auf der Halbkugel bleibt, für welche sie gelten, $\varphi'(\lambda)$ nicht unendlich wird und für ein kleines $\lambda$, für welches die erste Formel gilt, von der ersten Ordnung bleibt. Es findet in letzterer Beziehung nur für $\theta=\frac12\pi$ eine Ausnahme statt, in welchem Falle, da die beiden äusseren Intervalle verschwinden, $\varphi'(\lambda)$ überall den Ausdruck $\frac1\pi\cos\lambda$ erhält, welcher für kleine Werthe von $\lambda$ von der Ordnung Null ist, so dass also am Aequator eine unendlich grosse Dichtigkeit stattfindet.

Nach den für $\varphi'(\lambda)$ gefundenen Ausdrücken zerfällt das im allgemeinen Ausdrucke für $\rho$ enthaltene Integral in zwei andere, welche sich beziehungsweise von $0$ bis $\frac12\pi-\theta$ und von $\frac12\pi-\theta$ bis $\frac12\pi+\theta$ erstrecken, und deren erstes den einfachen Werth:
\[
 -4\cos\theta\sin\left(\frac14\pi-\frac12\theta\right)
\]
erhält. Das zweite, welches aus zwei Theilen besteht, ist, wenn man zunächst die Grenzen nicht berücksichtigt:
\[
 \frac{\sin\theta}{\pi}\int\frac{\sin\psi_1\cos\lambda}{\sin\frac12\lambda}\dd\lambda
 -\frac{2\cos\theta}{\pi}\int\psi_1\cos\frac12\lambda\dd\lambda,
\]
und erhält durch theilweise Integration des letzten Theils die Form:
\[
 \frac{\sin\theta}{\pi}\int\frac{\sin\psi_1\cos\lambda}{\sin\frac12\lambda}\dd\lambda
 +\frac{4\cos\theta}{\pi}\int\frac{\partial\psi_1}{\partial\lambda}\sin\frac12\lambda\dd\lambda
 -\frac4\pi\psi_1\cos\theta\sin\frac12\lambda.
\]
Bemerkt man, dass nach der $\psi_1$ bestimmenden Gleichung an den Grenzen resp. $\psi_1=\pi$ und $\psi_1=0$ ist, so erhält man beim Uebergange zu den Grenzen aus dem vom Integralzeichen freien Gliede einen Werth, welcher dem oben für das erste Integral gefundenen entgegengesetzt ist. Es bleiben daher nur die beiden ersten Glieder, welche durch Substitution der aus der Gleichung für $\psi_1$ sich ergebenden Ausdrücke:
\[
 \sin\psi_1=\frac{\sqrt{\sin^2\theta-\cos^2\lambda}}{\sin\theta\sin\lambda},
\qquad
 \frac{\partial\psi_1}{\partial\lambda}=-\frac{\cos\theta}{\sin\lambda}\frac1{\sqrt{\sin^2\theta-\cos^2\lambda}}
\]
die Form erhalten:
\[
 \frac1\pi\int\frac{\cos\lambda}{\sin\lambda\sin\frac12\lambda}\sqrt{\sin^2\theta-\cos^2\lambda}\dd\lambda
 -\frac{4\cos^2\theta}{\pi}\int\frac{\sin\frac12\lambda}{\sin\lambda}\frac{\dd\lambda}{\sqrt{\sin^2\theta-\cos^2\lambda}},
\]
in welcher die Integrale von $\lambda=\frac12\pi-\theta$ bis $\lambda=\frac12\pi+\theta$ zu nehmen sind. Führt man als neue Veränderliche
\[
 z=\frac{\cos\lambda}{\sin\theta}
\]
ein, so werden die Grenzen für diese $+1$ und $-1$, oder wenn man das Vorzeichen ändert, $-1$ und $+1$, und man erhält, wenn man in den Ausdruck für $\rho$ einsetzt und zugleich berücksichtigt, dass $\varphi(\pi)$ verschwindet:
\[
 \rho=\frac{1}{2\sqrt2\pi^2}\int_{-1}^{+1}\left(\left(\frac{2-z\sin\theta}{1-z^2\sin^2\theta}\right)\cos^2\theta-z\sin\theta\right)
 \frac{\dd z}{\sqrt{(1-z^2)(1-z\sin\theta)}}.
\]
Dieses Resultat enthält nur scheinbar einen von den elliptischen Integralen der dritten Gattung abhängigen Bestandtheil. Da nämlich die Grenzen zwei der aufeinander folgenden Werthe sind, für welche das Radical verschwindet, so wird der Ausdruck, auf die gewöhnliche Form gebracht, nur sogenannte vollständige elliptische Integrale enthalten\footnote{Fundamenta nova theoriae functionum ellipticarum, auctore C. G. J. Jacobi, p. 14. C. G. J. Jacobi's gesammelte Werke, Bd. I, S. 67. K.} und folglich nach einem schönen von Legendre herrührenden Satze auf die erste und zweite Gattung zurückgeführt werden können.

Zum Schluss wollen wir noch bemerken, wie der in der oben erwähnten Abhandlung behandelte Fall des Problems für eine Fläche, welche nur wenig von der Kugelfläche abweicht, ohne Reihenentwickelung auf den Fall der Kugel zurückgeführt werden kann.

Es sei $r=1+\gamma z$ die Gleichung der Fläche, worin $\gamma$ eine kleine Constante, deren höhere Potenzen vernachlässigt werden sollen, und $z$ eine Function von $\theta$ und $\varphi$ bezeichnet. Kommt man überein, den Werth, welchen eine Function von $\theta$, $\varphi$ erhält, wenn darin $\theta$, $\varphi$ in $\theta'$, $\varphi'$ verwandelt werden, durch einen hinzugefügten Accent zu bezeichnen, nennt $ds'$ das zu $\theta'$, $\varphi'$ gehörige Element der Fläche, und $p$ die Entfernung der den Coordinatenverbindung $\theta$, $\varphi$ und $\theta'$, $\varphi'$ entsprechenden Punkte derselben, so erfordert die Aufgabe, dass der Gleichung:
\[
 V=\int \rho'\frac{ds'}{p},
\]
worin sich die doppelte Integration über die ganze Fläche, oder von $\theta'=0$, $\varphi'=0$ bis $\theta'=\pi$, $\varphi'=2\pi$ erstreckt, für alle Werthe von $\theta$, $\varphi$ in demselben Umfange genüge. Ist $d\sigma'$ das dem Elemente $ds'$ entsprechende Element der Kugelfläche, welches von demselben Radienvectoren wie dieses ausgeschnitten wird, und $q$ die Entfernung der auf der Kugelfläche zu $\theta$, $\varphi$ und $\theta'$, $\varphi'$ gehörigen Punkte, so hat man sogleich durch die einfachsten geometrischen Betrachtungen:
\[
 ds'=(1+2\gamma z')d\sigma'
\]
und
\[
 p=q\left(1+\frac12\gamma(z+z')\right).
\]
Setzt man diese Ausdrücke ein, so wird unsere Gleichung:
\[
 V=\int \rho'\left(1+\gamma\left(\frac32z'-\frac12z\right)\right)\frac{d\sigma'}{q}
\]
oder:
\[
 V+\frac12\gamma z\int \rho'\frac{d\sigma'}{q}
 =\int \rho'\left(1+\frac32\gamma z'\right)\frac{d\sigma'}{q}.
\]
Da in Folge der vorletzten Gleichung das in der letzten mit $\gamma$ multiplicirte Integral von $V$ nur um eine Grösse der Ordnung $\gamma$ verschieden ist, so kann $V$ für dasselbe gesetzt werden, und man erhält die Gleichung:
\[
 \left(1+\frac12\gamma z\right)V=\int \rho'\left(1+\frac32\gamma z'\right)\frac{d\sigma'}{q},
\]
durch welche die Aufgabe auf den Fall der Kugelfläche zurückgeführt ist. Man sieht, dass man den gegebenen Potentialwerth, mit $1+\frac12\gamma z$ multiplicirt, als für die Kugelfläche geltend zu betrachten und dann die für diese Voraussetzung bestimmte Dichtigkeit durch $1+\frac32\gamma z$ zu dividiren hat.

\end{document}
